\section{Kesimpulan dan Saran Pengembangan}

\subsection{Kesimpulan}
Berdasarkan perancangan, implementasi, dan pengujian pada laporan ini, diperoleh kesimpulan berikut.

\begin{enumerate}
  \item \textbf{Rancangan memenuhi syarat spesifikasi.} \algname{} adalah \textit{block cipher}
    jaringan Feistel 16 putaran dengan blok dan kunci 128 bit. Substitusi dilakukan oleh \textit{S-box}
    8 bit, transposisi oleh permutasi bit 64 bit, dan operasi tambahan berupa penjumlahan modulo
    $2^{32}$ serta rotasi. Setiap putaran memakai kunci putaran yang dibangkitkan dari kunci utama oleh
    \textit{key schedule}.

  \item \textbf{Implementasi berfungsi pada kelima mode operasi.} ECB, CBC, CFB, OFB, dan CTR
    dilengkapi \textit{padding} PKCS\#7 dan verifikasi integritas \textit{encrypt-then-MAC}
    dengan CMAC. Pada seluruh pengujian, hasil dekripsi identik dengan berkas asli, sedangkan
    cipherteks yang diubah, kunci yang salah, dan kontainer yang terpotong ditolak.

  \item \textbf{Difusi cipher memenuhi target.} Avalanche effect rata-rata 49,99\% untuk perubahan
    1 bit plainteks maupun kunci, dan uji SAC tidak menemukan pasangan bit \textit{input} dan
    \textit{output} yang bias (0 dari 16.384 sel di luar $4\sigma$). Difusi penuh dicapai pada putaran
    ke-4, sehingga 16 putaran memberi margin sekitar empat kali lipat.

  \item \textbf{Keacakan cipherteks bergantung pada mode.} CBC, CFB, OFB, dan CTR menghasilkan
    cipherteks yang mendekati seragam (entropi 7,9998 bit/byte) pada ketiga \textit{input} uji,
    sedangkan ECB tetap membocorkan pola plainteks karena blok yang sama menghasilkan cipherteks
    yang sama. Hal ini sifat mode ECB dan bukan kelemahan cipher.

  \item \textbf{Ada keterbatasan yang perlu dicatat.}
    \begin{itemize}
      \item \textit{S-box} bijektif dan tanpa titik tetap, tetapi \textit{nonlinearity} 98 dan
        \textit{differential uniformity} 10 lebih lemah daripada \textit{S-box} AES (112 dan 4).
      \item Belum dilakukan analisis ketahanan terhadap kriptanalisis diferensial dan linear
        (misalnya penghitungan lintasan), dan \algname{} belum dikriptanalisis oleh pihak lain.
        Hasil uji statistik pada Bab~5 tidak menggantikan analisis tersebut.
      \item Konstanta rotasi dan pengali permutasi dipilih sebagai titik awal dan belum dicari
        secara sistematis.
    \end{itemize}

  \item \textbf{Inspirasi sigma-delta tercapai sebagian.} Rancangan fungsi $F$ terinspirasi dari
    modulator sigma-delta dengan dua gagasan: (a) \textit{state} yang dibawa dari satu \textit{word}
    ke \textit{word} berikutnya, dan (b) sifat \textit{lossy} pada kuantisasi yang diharapkan
    mempersulit penelusuran \textit{input} dari \textit{output}. Gagasan (a) tercapai, karena
    \textit{output} tiap langkah bergantung pada seluruh \textit{input} sebelumnya dalam satu
    \textit{pass}. Gagasan (b) belum tercapai: kuantisasi diganti \textit{S-box} bijektif agar dapat
    menyertakan kunci, sehingga $F$ pada implementasi ini masih dapat dibalik untuk kunci tertentu
    (100.000 dari 100.000 pasangan uji). Sifat \textit{lossy} karenanya tidak dapat dijadikan
    alasan keamanan, dan ketahanan yang teramati berasal dari kunci, \textit{S-box}, rotasi,
    permutasi, dan struktur Feistel. Gagasan ini dapat ditelaah kembali pada pengembangan
    selanjutnya.
\end{enumerate}

\subsection{Saran Pengembangan (Future Works)}
\begin{itemize}
  \item \textbf{Menelaah ulang kuantisasi \textit{lossy}.} Merancang langkah kuantisasi yang benar-benar
    membuang sebagian informasi (misalnya memotong sebagian bit \textit{output}) tanpa melemahkan
    difusi, lalu mengujinya kembali dengan avalanche dan SAC. Uji properti fungsi $F$
    pada Bab~5 perlu disesuaikan karena kondisinya akan berbalik.
  \item \textbf{Analisis kriptanalisis diferensial dan linear.} Menghitung jumlah \textit{S-box}
    aktif dan batas probabilitas lintasan diferensial dan linear pada $r$ putaran, sebagai
    pelengkap uji statistik.
  \item \textbf{Optimasi konstanta.} Mencari konstanta rotasi (7, 19, 13, 29) dan pengali permutasi
    secara sistematis berdasarkan avalanche per putaran, bukan sekadar titik awal.
  \item \textbf{Memperbaiki kualitas \textit{S-box}.} Memakai konstruksi aljabar atau pencarian
    yang lebih ekstensif agar \textit{nonlinearity} dan \textit{differential uniformity} mendekati
    \textit{S-box} AES.
  \item \textbf{Pengujian lebih luas.} Menambah uji keacakan statistik standar dan data uji yang lebih
    beragam, serta kriptanalisis oleh pihak lain sebelum \algname{} dapat dianggap layak dipakai di
    luar keperluan pembelajaran.
\end{itemize}
